% !TeX root = ../ADCS & GNC Documentation.tex
\chapter{Control System}

\section{State-Space Model}

\[
\dot{x} =
\begin{bmatrix}
	0 & \frac{1}{2}I \\
	0 & 0
\end{bmatrix} x +
\begin{bmatrix}
	0 \\
	J^{-1}
\end{bmatrix} u
\]

\section{Control Law}

The controller state vector is defined as
\[
x =
\begin{bmatrix}
	q_{e,v} \\
	\omega
\end{bmatrix}
\]
where \(q_{e,v}\) is the vector portion of the attitude error quaternion and \(\omega\) is the body-frame angular velocity error. This is used in the resulting control law:

\[
\tau = u = -Kx
\]

Where \(u\) is the desired 3-axis body torques of the satellite, and \(K\) is the gain matrix derived from the LQR control design.

\section{LQR Control Design}

The control system is designed using a Linear Quadratic Regulator (LQR) formulation.

The gain matrix $K$ is computed by solving the discrete-time algebraic Riccati equation, following discretization of the continuous-time system using a zero-order hold assumption.

The LQR design is based on qualitative tuning parameters, including:

\begin{itemize}
	\item Maximum attitude error
	\item Maximum angular rate
	\item Maximum actuator torque
\end{itemize}

These values are not strict physical limits, but rather representative scaling values used to define the relative importance of each state and input. As a result, controller performance is governed by the relative weighting between attitude error, angular velocity, and control effort, rather than the absolute numerical values chosen.

The relationship between these qualitative tuning parameters and the resulting controller gains is established through the construction of the weighting matrices.

\subsection{Weighting Matrices}

The cost function minimized by the LQR controller is of the form:

\[
J = \sum (x^T Q x + u^T R u)
\]

The matrices $Q$ and $R$ define the relative weighting of the state error and control effort in the cost function.

\begin{itemize}
	\item $Q$ penalizes deviations in the state vector, including attitude error and angular velocity
	\item $R$ penalizes control input effort (actuator torque)
\end{itemize}

The weighting matrices are constructed using the previously defined qualitative parameters for maximum expected attitude error, angular rate, and actuator torque. These values are used as scaling factors to normalize each term in the cost function.

Specifically, each state and input is weighted inversely proportional to the square of its corresponding characteristic value. This ensures that states with smaller acceptable magnitudes (such as attitude error) are penalized more heavily than those with larger acceptable magnitudes.

As with the overall LQR design, the absolute values used in constructing $Q$ and $R$ are not physically prescriptive. Instead, the resulting controller behavior is determined by the relative weighting between the different state components and control effort. Adjusting these parameters effectively shifts the balance between responsiveness, control authority, and smoothness of the system response.

\section{Gain Design}

\begin{itemize}
	\item Discretization via zero-order hold
	\item Gain computation done using the discrete Riccati equation
\end{itemize}

\section{Variable Gain Control}

An optional variable gain control scheme is implemented to improve fine-pointing performance.

This approach utilizes two separate gain matrices:

\begin{itemize}
	\item A nominal gain matrix for coarse pointing
	\item A secondary gain matrix with increased weighting on quaternion error for fine pointing near the target
\end{itemize}

When the spacecraft approaches the target attitude, the controller transitions from the nominal gain matrix to the fine-pointing gain matrix.

Although a smoothing transition is applied between the two gain sets, this switching behavior introduces a small transient response in the actuator commands. In practice, this manifests as a slight jerk in the reaction wheel torques.

\begin{figure}[h]
	\centering
	% Insert figure here
	\caption{Transient response during gain transition (reaction wheel torque)}
\end{figure}

\section{Integrator Augmentation}

A method for incorporating an integrator term into the controller was implemented to reduce steady-state error.

However, it was found that the baseline controller already achieves extremely small steady-state errors. In practice, the introduction of an integrator provided no meaningful improvement in performance.

Instead, the presence of step changes in the target attitude resulted in increased transient instability when the integrator term was active. These effects outweighed any marginal benefit in steady-state accuracy.

Given that the current steady-state errors are already negligible, and that the real hardware system is expected to exhibit greater limitations due to actuator and sensor imperfections, the integrator-based controller was not retained in the final implementation.

\section{Feedforward Terms}

Feedforward terms are included to:

\begin{itemize}
	\item Track time-varying reference attitudes
	\item Compensate for gyroscopic coupling
\end{itemize}

\section{Actuator Allocation}

Desired body torque is mapped to reaction wheel torques using a wheel configuration matrix.

The matrix $G \in \mathbb{R}^{3 \times n}$ defines the orientation of each reaction wheel in the spacecraft body frame. Each column of $G$ corresponds to the unit vector of a reaction wheel spin axis expressed in body coordinates:

\[
G =
\begin{bmatrix}
	| & | &        & | \\
	g_1 & g_2 & \cdots & g_n \\
	| & | &        & |
\end{bmatrix}
\]

where each $g_i \in \mathbb{R}^3$ is a unit vector representing the axis of the $i$-th reaction wheel.

The relationship between individual wheel torques $u_{rw}$ and the resulting body torque $\tau$ is given by:

\[
\tau = G u_{rw}
\]

Since the system is overactuated for the OreSat/SENTINEL bus for redundancy (e.g., four wheels for three axes), the inverse mapping is computed using the Moore-Penrose pseudoinverse:

\[
u_{rw} = G^\dagger \tau
\]

where $G^\dagger$ is the pseudoinverse of the configuration matrix.

This formulation provides a least-squares solution that minimizes control effort while achieving the desired body torque. Care must be taken to ensure consistent sign conventions between the definition of $G$, the actuator model, and the applied torque commands.

\section{Assumptions and Limitations}

\begin{itemize}
	\item Rigid body dynamics
	\item Small-angle approximation in estimation
	\item Simplified actuator saturation modeling
\end{itemize}

\section{Known Issues}

\begin{itemize}
	\item Sensitivity to quaternion sign conventions
	\item Reaction wheel allocation sign consistency
	\item Dependence on accurate frame transformations
	\item Small transient torque discontinuity during variable gain transitions
\end{itemize}